\label{der:8-2}
写$\adj(tI-A)=\sum_{j=0}^{n-1}B_jt^j$；每个余子式次数至多$n-1$，故这种有限展开存在。代入伴随恒等式并按$t$的次数比较：
\[B_{n-1}=I,\quad B_{j-1}-AB_j=c_jI\ (1\le j\le n-1),\quad -AB_0=c_0I.\]
从最高次往下回代，得到
\[B_0=A^{n-1}+c_{n-1}A^{n-2}+\cdots+c_1I.\]
再代入常数项$-AB_0=c_0I$，把右边移到左边，正好是$A^n+c_{n-1}A^{n-1}+\cdots+c_0I=0$。$n=1$时直接由$B_0=I,-A=c_0I$得到相同结论，毋须中间递推。

为统一求逆式，记$c_n=1$。若$c_0\ne0$，提取$A$得到
\[A\left(\sum_{j=1}^n c_jA^{j-1}\right)=-c_0I.\]
括号为$A$的多项式，与$A$可交换，所以左右乘都等于$-c_0I$；乘$-\frac{1}{c_0}$即构造出双侧逆。反向$c_0=(-1)^n\det A$，非零也正是可逆条件。

若$p=q\chi_A+r$，在矩阵多项式中，由有限和与幂的交换律，$(q\chi_A)(A)=q(A)\chi_A(A)=0$，所以$p(A)=r(A)$。这说明带余除法降次的每一步都有代数依据，不依赖$A$是否可对角化。
